\documentclass[compress,dvips,xcolor={dvipsnames},t]{beamer}

\usepackage[francais]{babel}
\usepackage{xspace}

\newcommand{\cpp}{\mbox{C \hspace*{-2.5mm} \raise 0.7mm\hbox{${\scriptscriptstyle ++}$}}\xspace}

\title{Présentation Personnelle}
\author{Christian Rinderknecht}
\date{16 janvier 2003}

\begin{document}

\frame{\maketitle}

% ------------------------------------------------------------------------
% Introduction

\begin{frame}
\frametitle{Parcours en raccourci}

\begin{itemize}
  \item Doctorat à l'INRIA au projet Cristal (1993-1997),
        sous le signe de l'application de techniques formelles à un
        langage télécom (réalisation d'un frontal de compilateur
        \mbox{ASN.1}, preuves de correction d'un codage/décodage).

  \item Depuis lors, alternance de postes académiques et industriels,
        avec comme point commun la modélisation de problèmes nouveaux,
        de solutions nouvelles et le développement de logiciels
        (ingénieur de recherche).

  \item Dernièrement (2001-2002), chercheur à l'\emph{Information and
        Communications University} (Corée du Sud).

\end{itemize}
\end{frame}


\begin{frame}
\frametitle{Les thématiques}

\begin{itemize}
  \item Une thématique d'application fut l'ingénierie des protocoles
        télécoms: spécification, méthodes et outils de production
        de tests. (trois ans)

  \item À Alcatel et PolySpace la thématique était le cycle de
        développement du logiciel: compilation, analyse statique de
        programmes, tests etc. (deux ans)

  \item Dans tous les cas, ma mission était de mener et d'appliquer
        une recherche, l'approche était rigoureuse (ou théorique),
        souvent guidée par l'aspect «~langage~» des problèmes.

\end{itemize}
\end{frame}

%% \begin{frame}
%% \frametitle{Expérience à PolySpace Technologies S.A.}

%% \begin{itemize}
%%   \item \emph{Start-up} INRIA (Rocquencourt).

%%   \item Développement et vente d'un logiciel d'analyse statique de
%%         code.

%%   \item Mes missions:

%%         \begin{itemize}

%%            \item maintenance et évolution de la base de test;

%%            \item conception et développement;

%%            \item études clients, commerce technique ponctuel.

%%         \end{itemize}
%% \end{itemize}

%% \end{frame}

\begin{frame}
\frametitle{Collaborations}

\begin{itemize}

  \item Durant le doctorat, collaboration étroite avec France Télécom
        R\&D (Lannion) --- condition propice à la création du pôle de
        compétence \mbox{ASN.1}.

  \item À l'Institut National des Télécommunications (INT), j'ai
        participé à deux projets \textsc{rnrt} (\textsc{castor} et
        \textsc{platonis}), avec pour partenaires Alcatel, France
        Télécom R\&D, Vérilog, le LaBRI, l'université de
        Clermont-Ferrand.

%%   \item À l'Institut National des Télécommunications, j'ai participé
%%         au projet \textsc{rnrt castor}: conception d'un atelier de
%%         services pour les plates-formes \textsc{corba} et Réseaux
%%         Intelligents (en collaboration avec Alcatel, France Télécom
%%         R\&D et Vérilog).

%%   \item Toujours à l'\textsc{int}, j'ai participé au projet
%%         \textsc{rnrt platonis} de plate-forme pour l'expérimentation
%%         et la validation de protocoles et de services sur des
%%         architectures mobiles (\textsc{wap}, \textsc{gprs}), en
%%         particulier l'étude des terminaux et les méthodes formelles
%%         pour le test d'interfonctionnement.

\end{itemize}

\end{frame}


\begin{frame}
\frametitle{Enseignement}

\begin{itemize}

  \item Monitorat (1993-1996) à l'université René Diderot (Paris VII)
        en IUP et DEUG: TD d'informatique (algorithmique, Pascal et
        \cpp{}), d'algèbre linéaire, de calcul intégral; conception et
        suivi des projets estudiantins, préparation des examens,
        soutenances.

  \item TD et quelques TP et cours à l'Institut Supérieur
        d'Électronique de Paris: programmation \cpp. (1993-1996)

  \item Colleur en informatique (algorithmique et programmation Caml)
        en Math. Spé. (Lycée Stanislas). (1996-1997)

  \item Formation continue Pascal (\textsc{ecf}, 1993) et \cpp{}
        (\textsc{Unilog Formation}, 1999)

  \item À l'INT: encadrement de projets de fin d'étude, d'un stagiaire
        de DEA, TP en formation initiale et continue à SDL
        (ponctuellement, 1998-2000)

\end{itemize}

\end{frame}


\begin{frame}
\frametitle{\emph{Visiting Professor at Information and Communications University} (2001-2002)}

\begin{itemize}

  \item Recherche à plein temps (\mbox{ASN.1}, test
        d'interfonctionnement) et étude de nouveaux sujets tels que
        XML et les réseaux mobiles \emph{ad hoc}.

        \emph{Abstract Syntax Notation One} (\mbox{ASN.1}) est un
	langage normalisé pour la définition de types dont les valeurs
	peuvent être échangées à travers un réseau entre deux
	applications communicantes, indépendament de leur possible
	hétérogénéité. Applications: gestion de réseaux, courriels
	sécurisés, téléphonie mobile, voix sur IP, communications
	sol-air, commerce électronique, et même bases de données
	(\textsc{bd}).

        J'ai résolu formellement l'important problème de la validation
        des spécifications \mbox{ASN.1} en utilisant une technique
        d'un domaine différent (les contraintes).

\end{itemize}

\end{frame}


\begin{frame}
\frametitle{Nouveaux axes de recherche}

Symbiose entre le \emph{web} et les \textsc{bd}: les sites \emph{web}
reposent de plus en plus sur des \textsc{bd}, les collections de pages
chaînées du \emph{web} sont une \textsc{bd} tentante.

Il faut réinventer une approche \textsc{bd} du \emph{web} car celui-ci
contient des données dites semi-structurées, c-à-d. auto-descriptives
et irrégulières. Domaines affines:

\begin{itemize}

  \item systèmes de types et analyse par contraintes des langages de
        requêtes,

  \item syntaxes de transfert (codages pour une faible bande passante),

  \item interfonctionnement de systèmes hétérogènes (p. ex. ASN.1-XML),

  \item réingénierie des sites \emph{web} (migration, XML
        \emph{wrappers}),

  \item protocoles pour les nouveaux terminaux mobiles (assistants
        numériques).

\end{itemize}

\end{frame}



\begin{frame}
\frametitle{Axes d'enseignement}

\begin{itemize}

  \item algorithmique, langages de programmation, compilation;

  \item environnements de développement Windows;

  \item environnement de développement Unix (\textsf{make},
        \textsf{cvs}, \textsf{emacs}, \textsf{sh}, \textsf{perl}
        etc.).

  \item \emph{middleware} (SOAP, .NET);

  \item XML;

  \item ingénierie des protocoles.

\end{itemize}

\end{frame}



\begin{frame}
\frametitle{Conclusion}

\begin{itemize}

  \item Formation doctorale à l'INRIA (recherche appliquée);

  \item enseignement à l'université et dans le privé;

  \item expérience dans l'ingénierie de recherche;

  \item publications des travaux réalisés.

\end{itemize}

\end{frame}

\end{document}

%%% Local Variables:
%%% mode: latex
%%% TeX-master: t
%%% End:
